\section{Câu 4.14}
The following contingency table summarizes supermarket transaction data, where \textit{hot dogs} refers to the transactions containing hot dogs, $\overline{hot\ dogs}$ refers to the transactions that do not contain hot dogs, \textit{hamburgers} refers to the transactions containing hamburgers, and $\overline{hamburgers}$ refers to the transactions that do not contain hamburgers.

\begin{center}
    \begin{tabular}{|c|c|c|c|}
        \hline
                                & $hot\ dogs$ & $\overline{hot\ dogs}$ & $\Sigma_{row}$ \\
        \hline
        hamburgers              & 2000        & 500                    & 2500           \\
        $\overline{hamburgers}$ & 1000        & 1500                   & 2500           \\
        \hline
        $\Sigma_{col}$          & 3000        & 2000                   & 5000           \\
        \hline
    \end{tabular}
\end{center}

\begin{itemize}
    \item[a.] Suppose that the association rule ``$hot\ dogs \Rightarrow hamburgers$'' is mined. Given a minimum support threshold of 25\% and a minimum confidence threshold of 50\%, is this association rule strong?
    \item[b.] Based on the given data, is the purchase of hot dogs independent of the purchase of hamburgers? If not, what kind of correlation relationship exists between the two?
    \item[c.] Compare the use of the all\_confidence, max\_confidence, Kulczynski, and cosine measures with lift and correlation on the given data.
\end{itemize}

\paragraph{Lời giải}
